% titel: Blatt 2, A7: Grundbegriffe der Rechnung in der Chemie
% skript: Mortimer Kap. 4.6 (Reaktionsgleichungen, Stöchiometrie)

% Teil I: Werkzeuge, kein Lösungsweg.

Reaktionsgleichungen beschreiben, welche Ausgangsstoffe (Edukte) zu welchen
Produkten umgesetzt werden und in welchen Mengenverhaeltnissen. Einige
Grundbegriffe müssen praezise sitzen.

\begin{werkzeugkasten}[Begriffe der chemischen Rechnung]
\begin{itemize}
  \item \textbf{Edukt / Produkt}: Stoffe auf der linken Seite einer
        Reaktionsgleichung sind die Edukte (Ausgangsstoffe); die rechte Seite
        enthält die Produkte (gebildete Stoffe).
  \item \textbf{Stochiometrische Koeffizienten}: die Zahlen vor den Formeln in
        einer ausbalancierten Gleichung; sie geben die kleinste ganze Zahl der
        Stoffmengenverhältnisse an (z.~B. $2\,\mathrm{H_2} +
        1\,\mathrm{O_2} \to 2\,\mathrm{H_2O}$).
  \item \textbf{$l$, $s$, $aq$, $g$}: Aggregatzustands-Symbole --- $l$
        flüssig (liquid), $s$ fest (solid), $g$ gasförmig (gas),
        $aq$ in Wasser gelöst (aquatisch/aqueous).
  \item \textbf{Begrenzender Reaktand}: der Edukt-Stoff, der zuerst
        aufgebraucht wird und die maximal bildbare Produktmenge bestimmt; alle
        anderen liegen im Ueberschuss vor.
  \item \textbf{tatsächliche vs.\ theoretische Ausbeute}: die theoretische
        Ausbeute ist die maximal mögliche (aus Stochiometrie berechnete)
        Menge; die tatsächliche ist die real gewonnene Menge (niedriger,
        wenn die Reaktion nicht vollständig laeuft, Nebenreaktionen
        auftreten oder Stoffe verloren gehen).
\end{itemize}
\textit{Skript:} Mortimer Kap. 4.6 (Reaktionsgleichungen, stoichiometrie),
Übung 2/7.
\end{werkzeugkasten}

\begin{achtung}[Darauf achten]
\begin{itemize}
  \item Koeffizienten beziehen sich auf die ganzen Formeln, nicht auf einzelne
        Elemente.
  \item ``aq'' meint in Wasser gelöst, nicht das Wasser selbst.
  \item Der begrenzende Reaktand bestimmt die Ausbeute; man muss ihn mit der
        molaren Stochiometrie (nicht mit der Masse) identifizieren.
\end{itemize}
\end{achtung}
